\paragraph{Prospective workflow feasibility pilot.} We prospectively fixed 12 telecom tasks and one run per task for each of two Haiku workflows (standard versus a verification instruction), with a common Haiku user simulator, randomized workflow order, a 40-step cap, and deterministic outcome verification. The fixed monetary stopping rule yielded 18/24 completed runs and 9/12 complete pairs; 12 completed runs reached the step cap. Standard and verification completed 9 and 9 runs, with 2 and 2 successes, respectively. Across all 12 planned tasks, the conservative missing-outcome bounds are [-0.417, 0.083] for net preference and [-0.250, 0.250] for success difference (verification minus standard). These are deterministic finite-cohort bounds, not sampling confidence intervals, and conservatively ignore partial one-arm information. The combined ledger accounts for \$3.9456 under a \$4.00 cap, including the retained \$0.0079 reservation from a quota-blocked initial model comparison that produced no completed trajectory. This small laboratory pilot establishes execution and measurement feasibility; limited success, truncation and incomplete pairs preclude a workflow-superiority conclusion. It does not involve live production users or production A/B randomization.
